% Класс документа с параметрами: кегль шрифта, шрифт с засечками для формул,
% соотношение сторон слайда
\documentclass[14pt,mathserif,aspectratio=43]{beamer}

% Кодировка файла
\usepackage[utf8]{inputenc}
% Переносы для русского языка
\usepackage[russian]{babel}

\usepackage{xcolor}
% Цвет логотипа лаборатории
\definecolor{lab-blue}{RGB}{30,97,182}
% Ищется в яндексе по запросу «травяной цвет»
\definecolor{grass}{RGB}{93,161,48}

\usepackage[T1]{fontenc}
% Шрифт для формул
\usepackage[bitstream-charter]{mathdesign}

% Математические символы
\usepackage{amsmath}

% Перечёркивание текста
\usepackage{ulem}

\usepackage{bbm}

% Основной цвета beamer'а (по умолчанию — синий)
\setbeamercolor{structure}{fg=black}

% Путь к картинкам
\graphicspath{{images/}}

% Нумерация слайдов
\setbeamertemplate{footline}{%
\raisebox{7pt}{\makebox[\paperwidth]{%
\hfill\makebox[10pt]{\scriptsize\textcolor{gray}{\insertframenumber~~~~}}}}}

% Начинать нумерацию слайдов не с титульного слайда
% \let\otp\titlepage
% \renewcommand{\titlepage}{\otp\addtocounter{framenumber}{-1}}

% Сноска без номера
\newcommand\articlenote[1]{%
  \begingroup%
  \renewcommand\thefootnote{}\footnote{#1}%
  \addtocounter{footnote}{-1}%
  \endgroup%
}

% Убрать навигацию beamer'а по умолчанию
\beamertemplatenavigationsymbolsempty

% Параметры для создания титульной страницы
\title{Обзор дескрипторов локальных особенностей \vspace{-1em}}
\author{Федор Морозов}
\usepackage{epstopdf}
\institute{

    Научный руководитель к.ф-м.н. Конушин А. С. \\
    \medskip
    Факультет ВМК МГУ им. М.В.\,Ломоносова\\
    Лаборатория КГиМ\\
    Семинар <<Компьютерное зрение>>

    \medskip

    \includegraphics[scale=0.9]{lab_logo.eps}

    \medskip
}

% Время
\date{\small{19 мая 2014 г.}}

\begin{document}

%-------------------------------------------------------------------------------
\begin{frame}[plain]
    \titlepage
\end{frame}

\begin{frame}{Задачи}
    \begin{itemize}
        \item Ознакомиться с аналогичными обзорами прошлых лет
        \pause
        \item Выбрать набор дескрипторов для тестирования
		\pause
		\item Выбрать задачу для оценки качества дескрипторов
		\pause
		\item Реализовать решение выбранной задачи и протестировать различные пары детекторов и дескрипторов	
    \end{itemize}
    
\end{frame}

%-------------------------------------------------------------------------------
\begin{frame}{Актуальность}

    \begin{itemize}
        \item Все доступные обзоры были написаны не позднее 2011 года.
        \pause
        \item В 2011 -- 2013 годах было разработано множество новых подходов, были доработаны подходы на основе машинного обучения.
		\pause
		\item Также в последние годы выросло число бинарных дескрипторов локальных особенностей.
		\pause
        \item Сравнение проводится для значительного числа современных алгоритмов.
		\pause
		\item Тестирование проводится на новой прикладной задаче.
    \end{itemize}
\end{frame}

\begin{frame}{Бинарные дескрипторы}
    \begin{itemize}
        \item Дескриптор представляет собой битовую строку фиксированной длины.
        \pause
        \item В результате дескриптор получается более компактным, чем вещественные.
        \pause
        \item Обычно можно использовать расстояние Хэмминга для сопоставления дескрипторов.
        \pause
        \item Возможно построение ускоренных алгоритмов сопоставления и поиска (например,  Multi-index hashing).
    \end{itemize}
\end{frame}

\begin{frame}{Области применения дескрипторов}
    \begin{itemize}
        \item Классификация и детектирование объектов

        \item Сопровождение объектов

        \item 3D-реконструкция

        \item Распознавание жестов
		
        \item Поиск изображений

		\item Сшивка панорам
    \end{itemize}
\end{frame}

\begin{frame}{Стандартные данные для тестирования}
	\begin{itemize}
        \item Mikolajczyk: набор изображений одной сцены с известными преобразованиями гомографии между кадрами.
        \pause
        \item Brown: база фрагментов изображений, взятых вокруг пар соответствующих ключевых точек.
        \pause
        \item opencv: множество поворотов одного изображения.
    \end{itemize}
\end{frame}

\begin{frame}{Другие обзоры}
    \begin{itemize}

    	\item 2005 --- Изображения сцены с известной гомографией (GLOH, SIFT, PCA-SIFT)
    	\pause
    	\item 2010 --- Поиск изображений (SIFT, CHoG, MPEG-7)
    	\pause
		\item 2010 --- Распознавание объектов	(SIFT, использование цвета)
		\pause
		\item 2011 --- Сопровождение объектов (SIFT, SURF, Random Forest)
		\pause
		\item 2011 --- Изображения сцены с известной гомографией (SIFT, DAISY)
		\pause
		\item 2011 --- Детектирование пешеходов (FREAK, SURF, BRISK)
	\end{itemize}    
\end{frame}

\begin{frame}{Рассмотренные дескрипторы}
    \begin{itemize}
    	\item Популярные вещественные

    	\begin{itemize}
    		\item SIFT

    		\item SURF

    	\end{itemize}

		\item Простые бинарные

		\begin{itemize}

			\item BRIEF

    		\item ORB

    		\item BRISK

			\item FREAK

			\item A-KAZE

		\end{itemize}

		\item Использующие бустинг

		\begin{itemize}

			\item BinBoost (128 бит)

    		\item BinBoost (256 бит)

    		\item BGM

		\end{itemize}
	\end{itemize}
\end{frame}

\begin{frame}{Разбор панорам}
    Дана серия из пяти изображений с камеры на автомобиле. \\
    Известно, что три или более из этих изображений сняты в одной локации (относятся к одной панораме). \\
    Задача --- определить лишние кадры в серии. \\
    
    \bigskip

    Оценка качества --- доля правильно классифицированных изображений. \\

	\bigskip

    \includegraphics[scale=0.25]{1_1.jpg}
    \includegraphics[scale=0.25]{1_2.jpg}
	\includegraphics[scale=0.25]{1_3.jpg}
	\includegraphics[scale=0.25]{1_4.jpg}
	\includegraphics[scale=0.25]{1_5.jpg}


\end{frame}

\begin{frame}{Формальная постановка задачи}
    Входные данные: \\

    \medskip

    	$\{I_1, \dots, I_5\} \in \mathbb{R}^{300 \times 300}$ --- набор изображений в оттенках серого.

		\pause
	Выходные данные: \\

		\medskip

		$\{x_1, \dots, x_5\} \in \mathbb\{0, 1\}$, где $x_i = 1$, если изображение входит в панораму, иначе $x_i = 0$.
		\pause
	Оценка качества: \\	

		\medskip

		$Score = \frac{1}{5}\sum_{i=1}^5{\mathbbm{1}(x_i = y_i)}$, где $y_i$ --- правильная метка для \textit{i-го} изображения.

\end{frame}

\begin{frame}{Алгоритм решения}
    \begin{enumerate}
    		\pause
			\item Найдем и опишем локальные особенности на всех изображениях.
			\pause
    		\item Для каждой пары изображений найдем соответствия между дескрипторами.
    		\pause
    		\item Для каждой пары изображений оценим модель гомографии и вычислим число удовлетворяющих ей соответствий.
    		\pause
    		\item Построим матрицу соответствия между изображениями.
    		\pause
    		\item Выберем набор изображений, наиболее вероятно составляющих панораму.

		\end{enumerate}
\end{frame}

\begin{frame}{Реализация}
    Для реализации использовалась библиотека opencv и язык C++. \\
    \pause
    \medskip

    Для дескрипторов BinBoost, BGM и A-KAZE были использованы авторские реализации. \\
    \pause
    \bigskip

    Тестирование проводилось на всех доступных данных (1000 серий для разбора). \\
\end{frame}

\begin{frame}{Детекторы}
    Дескрипторы были протестированы в паре с различными детекторами.
    \pause
    \bigskip

    Основным детектором был выбран SURF.
	\pause
	\bigskip
	\pause
	Для некоторых дескрипторов также использовался авторский детектор. \\
	Такие тесты отмечены символом *.

\end{frame}

\begin{frame}{Результаты сравнения}
	
	\begin{tabular}{ | l | c | }
    \hline
    BINBOOST-128 & 0.904 \\
    \hline
	BRIEF & 0.902 \\
	\hline
	SIFT* & 0.899 \\
	\hline
	BINBOOST-256 & 0.898 \\
	\hline
	BGM & 0.893 \\
	\hline
	SURF & 0.893 \\
	\hline
	ORB & 0.892 \\
	\hline
	BRISK & 0.888 \\
	\hline
	SIFT & 0.884 \\
	\hline
	AKAZE* & 0.863 \\
	\hline
	ORB* & 0.840 \\
	\hline
	FREAK & 0.788 \\
	\hline
	BRISK* & 0.656 \\
	\hline

    \end{tabular}
\end{frame}

\begin{frame}{Результаты сравнения}
	\includegraphics[width=\linewidth]{score.eps}
	\articlenote{* использование авторского детектора}
\end{frame}

\begin{frame}{Сравнение времени вычисления}
	\includegraphics[width=\linewidth]{time.eps}
\end{frame}

\begin{frame}{Скорость работы}
	\includegraphics[width=\linewidth]{speed.eps}
\end{frame}

\begin{frame}{Резюме}
	\begin{itemize}
        \item Современные двоичные дескрипторы представляют достойную альтернативу дескрипторам SIFT и SURF.
        \pause
        \item При этом многие из них требуют значительно меньше вычислительных ресурсов на вычисление и сопоставление.
        \pause
        \item Результат работы с локальными особенностями сильно зависит от выбранного детектора.
    \end{itemize}
\end{frame}

\begin{frame}{Дальнейшие планы}
	\begin{itemize}
        \item Реализация тестирования на данных Mikolajczyk и Brown.
        \pause
        \item Переработка тестирующей системы библиотеки opencv.
        \pause
        \item Тестирование дескрипторов в поблочных алгоритмах вычитания фона и 3D-реконструкции.
    \end{itemize}
\end{frame}

\end{document}

